% Шаг метода Розена: антиградиент выводит из допустимого множества,
% его проекция на активную грань — нет.
\begin{figure}[htbp]
\centering
\begin{tikzpicture}[scale=1.0, >={Stealth[length=2.2mm]}]

% допустимый многогранник
\fill[blue!8] (0,0) -- (5.2,0) -- (5.2,2.0) -- (3.0,3.4) -- (0,2.3) -- cycle;
\draw[blue!55!black, thick]
  (0,0) -- (5.2,0) -- (5.2,2.0) -- (3.0,3.4) -- (0,2.3) -- cycle;
\node[blue!55!black, font=\footnotesize] at (1.3,1.1) {$S$};

% активная грань (нижняя сторона) подсвечена
\draw[orange!75!black, very thick] (0,0) -- (5.2,0);
\node[orange!60!black, font=\scriptsize, anchor=north] at (4.3,-0.12)
  {активное ограничение $f_j\T x = g_j$};

% текущая точка на грани
\coordinate (xk) at (2.1,0);
\fill[black] (xk) circle (2pt);
\node[font=\footnotesize, anchor=south east] at ($(xk)+(-0.04,0.08)$) {$x_k$};

% антиградиент: вниз-вправо, выводит наружу
\coordinate (ag) at (3.25,-1.15);
\draw[gray!70!black, very thick, ->] (xk) -- (ag)
  node[below right, font=\footnotesize] {$-\nabla\varphi(x_k)$};

% проекция на грань
\coordinate (pr) at (3.25,0);
\draw[orange!75!black, very thick, ->] (xk) -- (pr)
  node[midway, above, font=\footnotesize, yshift=1pt]
  {$s_k = -P_k \nabla\varphi(x_k)$};

% пунктир: составляющая, которую отбрасывает проектор
\draw[gray, dashed] (ag) -- (pr);
\node[gray!60!black, font=\scriptsize, anchor=west] at (3.32,-0.62)
  {отброшено проектором};

% отметка прямого угла
\draw[gray] (3.08,0) -- (3.08,-0.17) -- (3.25,-0.17);

\end{tikzpicture}
\caption{Шаг метода Розена. Антиградиент в точке $x_k$ выводит за активное
ограничение, поэтому его проектируют на ядро матрицы активных ограничений:
составляющая, нарушающая ограничение, отбрасывается, и остаётся направление
$s_k$ вдоль грани. Если $s_k = 0$, движение вдоль грани исчерпано, и
оптимальность проверяется по знаку множителей $u$ в~\eqref{eq:rosen-5}.}
\label{fig:rosen-step}
\end{figure}
